% ============================================================
% Ejercicio 2
% ============================================================
\section{Ejercicio 2}

\begin{tcolorbox}[enunciado]
¿Cuáles de los siguientes problemas pertenecen a P, NP ó coNP? ¿Cuáles a coNP?

Justifica tu respuesta argumentando tu conjetura o dando la demostración de pertenencia a las clases correspondientes.

\textbf{Nota:} Considerar que al menos una demostración de pertenencia a NP se deberá realizar utilizando la noción de No Determinismo, y al menos una utilizando la definición alternativa (con la noción de verificación).
\end{tcolorbox}

\begin{itemize}[leftmargin=*]
    \item \textbf{DOM} = $\{ \langle G, K \rangle \mid G \text{ tiene un conjunto dominante de } K \text{ nodos} \}$ \\
    \textit{Un subconjunto de vértices $C$ de $G$ es un conjunto dominante de $G$, si todo otro vértice de $G$ es adyacente a algún vértice de $C$.}
    
    \subsection*{Solución}
    % Escribe tu respuesta aquí

    \vspace{0.5cm}
    \item \textbf{L} = $\{ \langle G \rangle \mid G \text{ es una gráfica no dirigida que NO contiene ciclos Eulerianos} \}$
    
    \subsection*{Solución}
    % Escribe tu respuesta aquí

    \vspace{0.5cm}
    \item \textbf{HAM\_MOD} = $\{ \langle G_1, G_2 \rangle \mid G_1 \text{ tiene un ciclo Hamiltoniano y } G_2 \text{ no tiene un ciclo Hamiltoniano} \}$
    
    \subsection*{Solución}
    % Escribe tu respuesta aquí

    \vspace{0.5cm}
    \item \textbf{CONECTIVIDAD} = $\{ \langle G \rangle \mid G \text{ es una gráfica conexa} \}$
    
    \subsection*{Solución}
    % Escribe tu respuesta aquí

    \vspace{0.5cm}
    \item \textbf{(*)} \textbf{CONJ\_INDEP} = $\{ \langle G, k \rangle \mid \text{la gráfica } G \text{ tiene un conjunto independiente de tamaño al menos } k \}$ \\
    \textit{* Un conjunto independiente de los vértices de una gráfica $G$ es un conjunto tal que ninguno de sus vértices es adyacente a otro.}
    
    \subsection*{Solución}
    % Escribe tu respuesta aquí

    \vspace{0.5cm}
    \item \textbf{(*)} \textbf{K-COLORACIÓN} = $\{ \langle G \rangle \mid G \text{ tiene una coloración de sus vértices, con } K \text{ colores} \}$ \\
    \textit{Investiga qué es una coloración en gráficas.}
    
    \subsection*{Solución}
    % Escribe tu respuesta aquí

    \vspace{0.5cm}
    \item \textbf{(*)} \textbf{EC-CUADRATICAS} = $\{ \langle a, b, c \rangle \mid a, b, c \in \mathbb{Z}^+, \exists x, y \in \mathbb{Z}^+ \text{ tal que } ax^2 + by = c \}$
    
    \subsection*{Solución}
    % Escribe tu respuesta aquí

\end{itemize}

\vspace{0.5cm}
\textit{* Opcional para equipos de 3 o menos integrantes. Equipos de 4 integrantes es obligatorio.}